\chapter{锁}
\label{CH:LOCK}

大多数内核，包括 xv6，都交错执行多个活动。交错的一个来源是多处理器硬件：具有多个独立执行的 CPU 的计算机，例如 xv6 的 RISC-V。这些多个 CPU 共享物理 RAM，xv6 利用这种共享来维护所有 CPU 读写内核数据结构。这种共享带来了一种可能性：一个 CPU 正在读取数据结构，而另一个 CPU 正在更新它的中途，甚至多个 CPU 同时更新同一数据；如果没有精心设计，这种并行访问可能会产生不正确的结果或破坏的数据结构。即使在单处理器上，内核也可能在多个线程之间切换 CPU，导致它们的执行被交错。最后，修改与某些可中断代码相同数据的设备中断处理程序，如果中断发生在恰好的错误时间，可能会损坏数据。
\indextext{并发}
指的是由于多处理器并行、线程切换或中断而导致多条指令流交错的情况。

内核充满了并发访问的数据。例如，两个 CPU 可以同时调用 {\tt kalloc}，从而并发地从空闲列表头部弹出。内核设计者喜欢允许大量并发，因为它可以通过并行性提高性能，并提高响应性。然而，因此内核设计者必须在尽管存在这种并发的情况下说服自己代码的正确性。有许多方法可以实现正确的代码，有些比其他更容易推理。旨在实现并发下正确性的策略和支持它们的抽象被称为
\indextext{并发控制}
技术。

Xv6 根据情况使用多种并发控制技术；还有更多可能的。本章专注于一种广泛使用的技术：
\indextext{锁}。
锁提供互斥，确保一次只有一个 CPU 可以持有锁。如果程序员将锁与每个共享数据项关联，并且代码在使用项目时始终持有关联的锁，那么一次只有一个 CPU 会使用该项目。在这种情况下，我们说锁保护数据项。虽然锁是一种易于理解的并发控制机制，但锁的缺点是它们可能限制性能，因为它们串行化并发操作。

本章的其余部分解释 xv6 为什么需要锁、xv6 如何实现它们以及如何使用它们。


% The following is most relevant to lock-free code:
% 
% A key observation will be that if you look at some code in
% xv6, you must ask yourself if concurrent code could change
% the intended behavior of the code by modifying data (or hardware resources)
% it depends on.
% You must keep in mind that the compiler may turn a
% single C statement into several machine instructions,
% and that those instructions may execute in a way that is
% interleaved with instructions executing on other CPUs.
% That is, you cannot assume that lines of C code
% on the page are executed atomically.
% Concurrency makes reasoning about correctness difficult.

%% 
\section{竞态}
%% 

\begin{figure}[t]
\center
\includegraphics[scale=0.8]{fig/smp.pdf}
\caption{简化的 SMP 架构}
\label{fig:smp}
\end{figure}

作为我们需要锁的一个例子，考虑两个具有已退出子进程的进程在两个不同的 CPU 上调用
{\tt wait} 系统调用。{\tt wait} 释放子进程的内存。因此在每个 CPU 上，内核将调用 {\tt kfree}
来释放子进程的内存页。内核分配器维护一个空闲页面的链表：\lstinline{kalloc()} \lineref{kernel/kalloc.c:/^kalloc/} 从列表中弹出一页内存，\lstinline{kfree()}
\lineref{kernel/kalloc.c:/^kfree/} 将一页压入列表。为了获得最佳性能，我们可能希望两个父进程的 {\tt kfree} 并行执行，而无需等待对方，但鉴于 xv6 的 {\tt kfree} 实现，这是不正确的。

Figure~\ref{fig:smp} 更详细地说明了设置：空闲页面链表位于两个 CPU 共享的内存中，两个 CPU 使用加载和存储指令操作列表。（实际上，处理器有缓存，但概念上多处理器系统表现得好像有单个共享内存。）
If there were no concurrent
requests, you might implement a list \lstinline{push} operation as
follows:
\begin{lstlisting}[numbers=left,firstnumber=1]
    struct element {
      int data;
      struct element *next;
    };
    
    struct element *list = 0;
    
    void
    push(int data)
    {
      struct element *l;
   
      l = malloc(sizeof *l);
      l->data = data;
      l->next = list; (*@\label{line:next}@*)
      list = l;  (*@\label{line:list}@*)
   }
\end{lstlisting}

\begin{figure}[t]
\center
\includegraphics[scale=0.5]{fig/race.pdf}
\caption{Example race}
\label{fig:race}
\end{figure}
This implementation is correct if executed in isolation.
However, the code is not correct if more than one
copy executes concurrently.
If two CPUs execute
\lstinline{push}
at the same time,
both might execute line~\ref{line:next} as shown in Fig~\ref{fig:smp},
before either executes line~\ref{line:list},  which results
in an incorrect outcome as illustrated by Figure~\ref{fig:race}.
There would then be two
list elements with
\lstinline{next}
set to the same former value of
\lstinline{list}.
When the two assignments to
\lstinline{list}
happen at line ~\ref{line:list},
the second one will overwrite the first;
the element involved in the first assignment
will be lost.

第 \ref{line:list} 行的丢失更新是
\indextext{竞态}
的一个例子。竞态是一种内存位置被并发访问的情况，且至少一次访问是写入。竞态通常是错误的迹象，要么是丢失更新（如果访问是写入），要么是读取不完全更新的数据结构。竞态的结果取决于编译器生成的机器代码、两个涉及的 CPU 的时序，以及它们的内存操作被内存系统排序的方式，这可能使竞态引起的错误难以重现和调试。例如，在调试
\lstinline{push}
时添加打印语句可能会改变执行的时序，足以使竞态消失。

避免竞态的常用方法是使用锁。锁确保
\indextext{互斥}，
以便一次只有一个 CPU 可以执行
\lstinline{push}
的敏感代码行；这使上述情况不可能发生。上述代码的正确加锁版本只添加了几行（以黄色高亮显示）：
\begin{lstlisting}[numbers=left,firstnumber=6]
   struct element *list = 0;
   (*@\hl{struct lock listlock;}@*)
    	
   void
   push(int data)
   {
     struct element *l;
     l = malloc(sizeof *l); (*@\label{line:malloc}@*)
     l->data = data;
   
     (*@\hl{acquire(\&listlock);} @*)
     l->next = list;     (*@\label{line:next1}@*)
     list = l;           (*@\label{line:list1}@*)
     (*@\hl{release(\&listlock)}; @*)
   }
\end{lstlisting}
\lstinline{acquire}
和
\lstinline{release}
之间的指令序列通常被称为
\indextext{临界区}。
锁被认为是在保护
\lstinline{list}。

当我们说锁保护数据时，我们真正的意思是锁保护适用于数据的一组不变量。不变量是跨操作维护的数据结构的属性。通常，操作的正确行为取决于操作开始时不变量为真。操作可能暂时违反不变量，但必须在完成之前重新建立它们。例如，在链表情况下，不变量是
\lstinline{list}
指向列表中的第一个元素，且每个元素的
\lstinline{next}
字段指向下一个元素。
\lstinline{push}
的实现暂时违反此不变量：在第 \ref{line:next1} 行，
\lstinline{l}
指向
列表的下一个元素，但
\lstinline{list}
尚未指向
\lstinline{l}
（在第 \ref{line:list1} 行重新建立）。我们上面检查的竞态发生是因为第二个 CPU 执行了依赖于列表不变量的代码，而它们被（暂时）违反了。正确使用锁确保一次只有一个 CPU 可以在临界区操作数据结构，因此没有 CPU 会在数据结构不变量不成立时执行数据结构操作。

你可以将锁视为
\indextext{串行化}
并发临界区，使它们一次运行一个，从而保留不变量（假设临界区在隔离情况下是正确的）。你也可以将由相同锁保护的临界区视为彼此原子的，以便每个只看到来自早期临界区的完整更改集，而永远不会看到部分完成的更新。

Though useful for correctness, locks inherently
limit performance.  For example, if two processes call {\tt kfree}
concurrently, the locks will serialize the two critical sections,
so that there is no
benefit from running them on different CPUs.  We say that multiple
processes \indextext{conflict} if they want the same lock at the same
time, or that the lock experiences \indextext{contention}.  A major
challenge in kernel design is avoidance of lock contention
in pursuit of parallelism.  Xv6 does
little of that, but sophisticated kernels organize
data structures and algorithms specifically to avoid lock contention.  In the list
example, a kernel may maintain a separate free list per CPU and only touch
another CPU's free list if the current CPU's list is empty and it must steal
memory from another CPU.  Other use cases may require more
complicated designs.

The placement of locks is also important for performance.
For example, it would be correct to move
\lstinline{acquire}
earlier in
\lstinline{push},
before line~\ref{line:malloc}.
But this would likely reduce performance because then the calls
to
\lstinline{malloc}
would be serialized.
The section ``Using locks'' below provides some guidelines for where to insert
\lstinline{acquire}
and
\lstinline{release}
invocations.
%% 
\section{代码：锁}
%% 

请阅读 {\tt kernel/spinlock.h} 和 {\tt kernel/spinlock.c}。

Xv6 有两种类型的锁：自旋锁和睡眠锁。我们从自旋锁开始。Xv6 将自旋锁表示为
\indexcode{struct spinlock}
\lineref{kernel/spinlock.h:/struct.spinlock/}。
结构中的重要字段是
\lstinline{locked}，
当锁可用时该字段为零，当锁被持有时为非零。从逻辑上讲，xv6 应该通过执行如下代码获取锁：
\begin{lstlisting}[numbers=left,firstnumber=21]
   void
   acquire(struct spinlock *lk) // does not work!
   {
     for(;;) {
       if(lk->locked == 0) {  (*@\label{line:testlocked}@*)
         lk->locked = 1;      (*@\label{line:assign}@*)
         break;
       }
     }
   }
\end{lstlisting}
不幸的是，这种实现不能
保证多处理器上的互斥。
可能发生两个 CPU 同时
到达第 \ref{line:testlocked} 行，看到
\lstinline{lk->locked}
为零，然后都通过执行第 \ref{line:assign} 行获取锁。
此时，两个不同的 CPU 持有锁，
这违反了互斥属性。
我们需要一种方法来
使第 \ref{line:testlocked} 和 \ref{line:assign} 行作为一个
\indextext{原子}
（即不可分割的）步骤执行。

因为锁被广泛使用，
多核处理器通常提供一条指令，
可用于使
第 \ref{line:testlocked} 和 \ref{line:assign} 行
原子化。
在 RISC-V 上，这条指令是
\lstinline{amoswap r, a}。
\lstinline{amoswap}
读取内存地址 {\tt a} 处的值，
将寄存器 {\tt r} 的内容写入该地址，
并将其读取的值放入 {\tt r}。
也就是说，它交换寄存器和被寻址
内存位置的内容。
它使用特殊
硬件原子地执行此序列，以防止任何
其他 CPU 在读取和写入之间使用该内存地址。

可移植的 C 库调用
\lstinline{__sync_lock_test_and_set(addr, value)}
归结为
\lstinline{amoswap}
指令；
该函数返回 {\tt *addr} 的旧（交换的）内容。
这是在 {\tt acquire} 中编写循环的好方法：

\begin{lstlisting}[numbers=left,firstnumber=21]
  while(__sync_lock_test_and_set(&lk->locked, 1) != 0)
    ;
\end{lstlisting}

Xv6 的
\indexcode{acquire}
\lineref{kernel/spinlock.c:/^acquire/}
使用上述循环。
每次迭代将 1 交换到
\lstinline{lk->locked}
并检查先前的值；
如果先前的值为零，则我们已获取
锁，交换将设置
\lstinline{lk->locked}
为 1。
如果先前的值为 1，则其他某个 CPU
持有锁，我们将 1 原子地交换到
\lstinline{lk->locked}
的事实不会改变其值。

一旦获取锁，
\lstinline{acquire}
记录获取锁的 CPU
以供调试。
\lstinline{lk->cpu}
字段受锁保护
且只能在持有锁时更改。

函数
\indexcode{release}
\lineref{kernel/spinlock.c:/^release/}
与
\lstinline{acquire}
相反：
它清除
\lstinline{lk->cpu}
字段
然后释放锁。
从概念上讲，释放只需要将零赋给
\lstinline{lk->locked}。
C 标准允许编译器用多个存储指令实现赋值，
因此 C 赋值可能相对于
并发代码是非原子的。
Instead,
\lstinline{release}
uses the C library function
\lstinline{__sync_lock_release}
that performs an atomic assignment.
该函数也归结为 RISC-V
\lstinline{amoswap}
指令。
%% 
\section{代码：使用锁}
%% 
Xv6 在许多地方使用锁来避免竞态。
如上所述，
\lstinline{kalloc}
\lineref{kernel/kalloc.c:/^kalloc/}
和
\lstinline{kfree}
\lineref{kernel/kalloc.c:/^kfree/}
是一个很好的例子。
尝试练习 1 和 2 看看如果这些
函数省略锁会发生什么。
你可能会发现很难触发不正确的
行为，这表明很难可靠地测试代码
是否没有锁错误和竞态。
Xv6 可能还有尚未发现的竞态。

使用锁的一个难点是决定使用多少锁以及每个锁应该保护哪些数据和不变量。
有一些基本原则。
首先，任何时候一个 CPU 写入变量
而另一个 CPU 同时读取或写入它，
应该使用锁来防止两个
操作重叠。
其次，记住锁保护不变量：
如果不变量涉及多个内存位置，
通常它们都需要被单个锁保护
以确保不变量被维护。

上面的规则说明了何时需要锁，但没有说明何时不需要锁，为了效率不过度加锁很重要，
因为锁减少并行性。
如果并行性不重要，那么可以
安排只有一个线程而不用担心锁。
一个简单的内核可以通过拥有一个锁来在多处理器上做到这一点；
内核每次从用户空间进入内核时获取锁，
用于系统调用或中断；
内核在返回用户空间时释放锁。
许多单处理器操作系统已被转换以
使用这种方法在多处理器上运行，有时称为``大
内核锁''，但这种方法牺牲了并行性：一次只有一个
CPU 可以在内核中执行。
如果内核消耗
大量 CPU 时间，可以通过用不同的锁保护不同的对象或模块来获得更多并行性，
以便不同的 CPU 可以同时
在内核的不同部分执行。

作为粗粒度锁的一个例子，xv6 的 \lstinline{kalloc.c}
分配器有一个由单个锁保护的单个空闲列表。如果
不同 CPU 上的多个进程同时尝试分配页面，
每个都必须在 {\tt acquire} 中自旋等待轮次。
自旋浪费 CPU 时间，因为这不是有用的工作。
如果锁争用浪费了很大比例的 CPU 时间，
也许可以通过修改分配器
为每个 CPU 提供单独的空闲列表，每个列表有自己的锁，以允许真正
并行的分配来提高性能。

作为细粒度锁的一个例子，xv6 为每个文件有一个单独的锁，以便操作不同文件的进程通常可以
继续而不必等待对方的锁。如果希望允许
进程同时写入同一文件的不同区域，文件锁定
方案可以做得更细粒度。
最终锁粒度决策需要由性能
测量以及复杂性考虑来驱动。

随着后续章节解释 xv6 的每个部分，它们
会提到 xv6 使用锁
处理并发的例子。
作为预览，
Figure~\ref{fig:locktable}
列出了 xv6 中的所有锁。

\begin{figure}[t]
\center
\begin{tabular}{ll}
{\bf Lock} & {\bf Description} \\
\midrule
bcache.lock & Protects allocation of block buffer cache entries \\
cons.lock & Serializes read processing of console input \\
tx\_lock & Serializes access to console (uart) output hardware \\
ftable.lock & Serializes allocation of a struct file in file table \\
itable.lock & Protects allocation of in-memory inode entries \\
vdisk\_lock & Serializes access to disk hardware and queue of DMA descriptors \\
kmem.lock & Serializes allocation of memory \\
log.lock & Serializes operations on the transaction log \\
pipe's pi->lock & Serializes operations on each pipe \\
pid\_lock & Serializes increments of next\_pid \\
proc's p->lock & Serializes changes to process's state \\
wait\_lock & Helps wait avoid lost wakeups \\
tickslock & Serializes operations on the ticks counter \\
inode's ip->lock & Serializes operations on each inode and its content \\
buf's b->lock & Serializes operations on each block buffer \\
\end{tabular}
\caption{xv6 中的锁}
\label{fig:locktable}
\end{figure}

%% 
\newpage
\section{死锁和锁排序}
%% 

假设在 CPU C1 和 C2 上运行的函数都有一个点，每个都需要持有锁 A 和锁 B，并且它们以不同的顺序获取它们：

\begin{center}
  \input{fig/order.tex}
\end{center}

运气不好的话，C1 和 C2 可能都在完全相同的时刻执行它们的第一次获取；两者都可能成功，因为它们要求的是不同的锁。但然后 C1 和 C2 都必须在它们对 {\tt acquire()} 的第二次调用中等待，因为两个锁都已被另一个 CPU 持有。因为两个 CPU 都在等待对方，所以谁都不会释放锁，两者都将永远等待。这种情况被称为
\indextext{死锁}。

C1/C2 示例中的关键问题是两个 CPU 以不同的顺序获取锁。如果它们都先尝试获取 A，一个会先获取 A 然后 B 然后释放它们，然后另一个 CPU 可以继续。更一般地说，锁代码必须遵循以下规则以避免死锁：所有持有多个锁的代码路径必须以相同的顺序获取锁。这种全局锁获取顺序的需要意味着锁实际上是每个函数规范的一部分：调用者必须以导致锁按约定顺序获取的方式调用函数。

Xv6 有许多长度为二的锁顺序链，涉及每个进程的锁（每个
\lstinline{struct proc}
中的锁）由于
\lstinline{sleep}
的工作方式（参见 Chapter~\ref{CH:SLEEP}）。例如，
\lstinline{consoleintr}
\lineref{kernel/console.c:/^consoleintr/}
是处理输入字符的中断例程。当新行到达时，任何等待控制台输入的进程都应该被唤醒。为此，
\lstinline{consoleintr}
持有
\lstinline{cons.lock}
同时调用
\indexcode{wakeup}，后者获取等待进程的锁以唤醒它。因此，全局避免死锁的锁顺序包括规则
\lstinline{cons.lock}
必须在任何进程锁之前获取。文件系统代码包含 xv6 最长的锁链。例如，创建文件需要同时持有目录上的锁、新文件 inode 上的锁、磁盘块缓冲区上的锁、磁盘驱动程序的 \lstinline{vdisk_lock} 以及调用进程的 \lstinline{p->lock}。为了避免死锁，文件系统代码始终按上一句提到的顺序获取锁。

遵守全局避免死锁的顺序可能出人意料地困难。有时锁顺序与逻辑程序结构冲突，例如，也许代码模块 M1 调用模块 M2，但锁顺序要求先获取 M2 中的锁再获取 M1 中的锁。有时锁的标识事先不知道，也许是因为必须先持有一个锁才能发现要获取的下一个锁的标识。这种情况在文件系统中出现，因为它查找路径名中的连续组件，以及在 {\tt wait} 和 {\tt exit} 系统调用的代码中，因为它们搜索进程表寻找子进程。最后，死锁的危险通常是锁定方案可以做得多细粒度的约束，因为更多的锁通常意味着更多的死锁机会。避免死锁的需要通常是内核实现中的一个主要因素。

有时会出现一个问题是，如果 CPU 尝试获取同一 CPU 已持有的锁应该发生什么。一条推理线是这应该被允许：没有其他 CPU 可以持有该锁，所以不必担心 CPU 干扰彼此使用受保护的数据。允许 CPU 重新获取它已持有的锁的锁系统称为{\it 可重入}或{\it 递归}。另一方面，如果锁已被持有，即使在同一 CPU 上，也意味着操作可能已暂时违反且尚未恢复某些不变量；允许在不变量不成立时开始新操作似乎是错误的邀请。Xv6 持后一种观点，并禁止当前持有锁的 CPU 重新获取它。检测这种情况是 {\tt acquire} 中调用 {\tt holding} 的目的。

\section{锁和中断}
\label{s:lockinter}

%% 
Some xv6 spinlocks protect data that is used by
both threads and interrupt handlers.
For example, the
\lstinline{clockintr}
timer interrupt handler might increment
\indexcode{ticks} 
\lineref{kernel/trap.c:/^clockintr/}
at about the same time that a kernel
thread reads
\lstinline{ticks} 
in
\indexcode{sys_pause}
\lineref{kernel/sysproc.c:/ticks0.=.ticks/}.
The lock
\indexcode{tickslock}
serializes the two accesses.

The interaction of spinlocks and interrupts raises a potential danger.
Suppose
\indexcode{sys_pause}
holds
\indexcode{tickslock},
and its CPU is interrupted by a timer interrupt.
\lstinline{clockintr}
would try to acquire
\lstinline{tickslock},
see it was held, and wait for it to be released.
In this situation,
\lstinline{tickslock}
will never be released: only
\lstinline{sys_pause}
can release it, but
\lstinline{sys_pause}
will not continue running until
\lstinline{clockintr}
returns.
So the CPU will deadlock, and any code
that needs either lock will also freeze.

To avoid this situation, if a spinlock is used by an interrupt handler,
a CPU must never hold that lock with interrupts enabled.
Xv6 is more conservative: when a CPU acquires any
lock, xv6 always disables interrupts on that CPU.
Interrupts may still occur on other CPUs, so 
an interrupt's
\lstinline{acquire}
can wait for a thread to release a spinlock; just not on the same CPU.

Xv6 re-enables interrupts when a CPU holds no spinlocks; it must
do a little book-keeping to cope with nested critical sections.
\lstinline{acquire}
calls
\indexcode{push_off}
\lineref{kernel/spinlock.c:/^push_off/}
and
\lstinline{release}
calls
\indexcode{pop_off}
\lineref{kernel/spinlock.c:/^pop_off/}
来跟踪当前 CPU 上锁的嵌套级别。
当该计数达到零时，
\lstinline{pop_off}
恢复在最外层临界区开始时存在的中断使能状态。
\lstinline{intr_off}
和
\lstinline{intr_on}
函数分别执行 RISC-V 指令以禁用和启用中断。

重要的是
\indexcode{acquire}
严格在设置
\lstinline{lk->locked}
\lineref{kernel/spinlock.c:/sync_lock_test_and_set/}
之前调用
\lstinline{push_off}。
如果两者顺序颠倒，会有一个短暂的窗口期，锁被持有但中断是启用的，一个不幸时序的中断会使系统死锁。
类似地，重要的是
\indexcode{release}
只在
释放锁之后
\lineref{kernel/spinlock.c:/sync_lock_release/}
调用
\indexcode{pop_off}。

%% 
\section{指令和内存排序}
%% 

将程序视为按源代码语句出现的顺序执行是很自然的。对于单线程代码，这是一个合理的心智模型，但当多个线程通过共享内存交互时，这是不正确的。一个原因是编译器以不同于源代码暗示的顺序发出加载和存储指令，并且可能完全省略它们（例如通过在寄存器中缓存数据）。另一个原因是 CPU 可能为了性能而乱序执行指令。例如，CPU 可能注意到在串行指令序列中 A 和 B 彼此不依赖。CPU 可能先开始指令 B，要么因为它的输入比 A 的输入先准备好，要么为了重叠执行 A 和 B。

As an example of what could go wrong,
in this code for
\lstinline{push},
it would be a disaster if the compiler or CPU moved the
store corresponding to
line~\ref{line:next2} to a point after the
\lstinline{release}
on line~\ref{line:release}:
\begin{lstlisting}[numbers=left,firstnumber=1]
      l = malloc(sizeof *l);
      l->data = data;
      acquire(&listlock);
      l->next = list;   (*@\label{line:next2}@*)
      list = l;      
      release(&listlock);  (*@\label{line:release}@*)
\end{lstlisting}
If such a re-ordering occurred, there would be a window during
which another CPU could acquire the lock and
observe the updated
\lstinline{list},
but see an uninitialized
\lstinline{list->next}.

The good news is that compilers and CPUs help concurrent programmers
by following a set of rules called the \indextext{memory model}, and
by providing some primitives to help programmers control re-ordering.

To tell the hardware and compiler not to re-order,
xv6 uses
\lstinline{__sync_synchronize()} 
in both
\lstinline{acquire} \lineref{kernel/spinlock.c:/^acquire/}
and
\lstinline{release} \lineref{kernel/spinlock.c:/^release/}.
\lstinline{__sync_synchronize()}
is a \indextext{memory barrier}:
it tells the compiler and CPU to not reorder loads or stores across the
barrier.
The barriers in xv6's 
\lstinline{acquire}
and
\lstinline{release}
force order in almost all cases where it matters,
since xv6 uses locks around accesses to shared data.
Chapter~\ref{CH:LOCK2} discusses a few exceptions.

%% 
\section{Sleep locks}
%% 

% this material refers a lot to sleep, yield, etc, so maybe it should
% be later, after sleep/wakeup in sched.tex.
% maybe in lock2.tex

Sometimes xv6 needs to hold a lock for a long time. For example, the
file system (Chapter~\ref{CH:FS}) keeps a file locked while reading
and writing its content on the disk, and these disk operations can
take tens of milliseconds. Holding a spinlock that long would lead to
waste if another process wanted to acquire it, since the acquiring
process would waste CPU for a long time while spinning. Another
drawback of spinlocks is that a process cannot yield the CPU while
retaining a spinlock; we'd like to do this so that other processes can
use the CPU while the process with the lock waits for the disk.
Yielding while holding a spinlock is illegal because it might
lead to deadlock if a second thread then tried to acquire the spinlock;
since {\tt acquire} doesn't yield the CPU, the second thread's
spinning might
prevent the first thread from running and releasing the lock.
% it's true that timer interrupts might force a switch from the
% second back to the first thread -- but not if the second
% already holds one or more locks.
% it's also true that the danger is most obvious on a uniprocessor,
% but it can also arise with a longer chain of threads on
% a multiprocessor.
Yielding while holding a lock would also violate
the requirement that interrupts must be off while a spinlock is held.
Thus we'd like a type of lock that yields the CPU while waiting to
acquire, and allows yields (and interrupts) while the lock is held.

Xv6 provides such locks in the form of
\indextext{sleep-locks}.
\lstinline{acquiresleep}
\lineref{kernel/sleeplock.c:/^acquiresleep/}
yields the CPU while waiting,
using techniques that will be explained in
Chapter~\ref{CH:SLEEP}.
At a high level, a sleep-lock has a
\lstinline{locked}
field that is protected by a spinlock, and 
\lstinline{acquiresleep} 's
call to
\lstinline{sleep}
atomically yields the CPU and releases the spinlock.
The result is that other threads can execute while
\lstinline{acquiresleep}
waits.

Because sleep-locks leave interrupts enabled, they cannot be
used in interrupt handlers.
Because
\lstinline{acquiresleep}
may yield the CPU,
sleep-locks cannot be used inside spinlock critical
sections (though spinlocks can be used inside sleep-lock
critical sections).

Spin-locks are best suited to short critical sections,
since waiting for them wastes CPU time;
sleep-locks work well for lengthy operations.

%% 
\section{Real world}
%% 
Programming with locks remains challenging despite years of research
into concurrency primitives and parallelism.
It is often best to conceal locks within 
higher-level constructs like synchronized queues, although xv6 does not
do this.  If you program with locks, it is wise to use a tool that attempts to
identify races, because it is easy to miss an invariant that requires
a lock.

Most operating systems support POSIX threads (Pthreads), which allow a
user process to have several threads running concurrently on different
CPUs.  Pthreads has support for user-level locks, barriers, etc.
Pthreads also allows a programmer to optionally specify that a lock
should be re-entrant.

Supporting Pthreads at user level requires support from the operating
system. For example, it should be the case that if one pthread blocks
in a system call, another pthread of the same process should be able
to run on that CPU.  As another example, if a pthread changes its
process's address space (e.g., maps or unmaps memory), the kernel must
arrange that other CPUs that run threads of the same process update
their hardware page tables to reflect the change in the address space.

It is possible to implement locks without atomic
instructions~\cite{lamport:bakery}, but it is expensive, so most
operating systems use atomic instructions.

Locks can be expensive if many CPUs try to acquire the same lock
at the same time.  If one CPU has a lock
cached in its local cache, and another CPU must acquire the lock, then the
atomic instruction to update the cache line that holds the lock must move the line
from the one CPU's cache to the other CPU's cache, and perhaps
invalidate any other copies of the cache line.  Fetching a cache line from
another CPU's cache can be orders of magnitude more expensive than
fetching a line from a local cache.

To avoid the expenses associated with locks, many operating systems
use lock-free data structures and algorithms~\cite{herlihy:art,mckenney:rcuusage}.
For example, it is possible to implement a linked list like the one in
the beginning of the chapter that requires no locks during list
searches, and one atomic instruction to insert an item in a list.
Lock-free programming is more complicated, however, than programming
locks; for example, one must worry about instruction and memory
reordering.  Programming with locks is already hard, so xv6 avoids the
additional complexity of lock-free programming.

%% 
\section{Exercises}
%% 

\begin{enumerate}
  
\item Comment out the calls to
\lstinline{acquire}
and
\lstinline{release}
in
\lstinline{kalloc}
\lineref{kernel/kalloc.c:/^kalloc/}.
This seems like it should cause problems for
kernel code that calls
\lstinline{kalloc};
what symptoms do you expect to see?
When you run xv6, do you see these symptoms?
How about when running
\lstinline{usertests}?
If you don't see a problem, why not?
See if you can provoke a problem by inserting
dummy loops into the critical section of
\lstinline{kalloc}.

\item Suppose that you instead commented out the
locking in
\lstinline{kfree} 
(after restoring locking in
\lstinline{kalloc}).
What might now go wrong? Is lack of locks in
\lstinline{kfree}
less harmful than in
\lstinline{kalloc}?

\item If two CPUs call
\lstinline{kalloc}
at the same time, one will have to wait for the other,
which is bad for performance.
Modify 
\lstinline{kalloc.c}
to have more parallelism, so that simultaneous
calls to
\lstinline{kalloc}
from different CPUs can proceed without waiting for each other.

\item Write a parallel program using POSIX threads, which is supported on most
operating systems. For example, implement a parallel hash table and measure if
the number of puts/gets scales with increasing number of CPUs.

\item Implement a subset of Pthreads in xv6.  That is, implement a user-level
thread library so that a user process can have more than 1 thread and arrange
that these threads can run in parallel on different CPUs.  Come up with a
design that correctly handles a thread making a blocking system call and
changing its shared address space.

\end{enumerate}

% LocalWords:  CPUs uniprocessor interruptible allocator kalloc kfree
% LocalWords:  struct malloc sizeof listlock invariants spinlock lk
% LocalWords:  amoswap cpu bcache ftable itable inode vdisk DMA kmem
% LocalWords:  pid proc's tickslock inode's ip buf's proc SCHED sys
% LocalWords:  consoleintr wakeup clockintr intr acquiresleep POSIX
% LocalWords:  Pthreads pthread unmaps TLB IPIs SMP firstnumber int
% LocalWords:  structure's spinlocks addr tx uart pipe's wakeups A's
% LocalWords:  usertests
